\documentclass[12pt]{article}

\usepackage[
margin=0.4in,
bottom=0.55in,
top=0.4in
]{geometry}


\usepackage{array}
\usepackage{enumitem}
\usepackage{multicol}

\setlist[itemize]{
	topsep=0pt,
	partopsep=0pt,
	parsep=0pt,
	itemsep=0pt
}

\newcolumntype{L}[1]{>{\raggedright\arraybackslash}p{#1}}
\newcommand{\checkbox}{\fbox{\rule{0pt}{0.8ex}\rule{0.8ex}{0pt}}}

\usepackage{xcolor}
\definecolor{solutionred}{HTML}{7F1D1D}
\newcommand{\sol}[1]{{\color{solutionred}\small\strut #1\strut}}
\newcommand{\answerline}[1]{\makebox[0pt][l]{\raisebox{0.6ex}{\sol{#1}}}}
\begin{document}
	
	\begin{center}
		{\Large\bfseries C3--C4: ESCUCHAR A LOS USUARIOS}
		
		\vspace{0.1cm}
		
		Emprendimiento e Innovación Grado 10\\[-1mm]{\scriptsize\color{solutionred}Modelo pedagógico: toda evidencia y respuesta es simulada; no son datos del colegio.}
		
		\vspace{0.3cm}
		
		\begin{tabular}{L{0.12\textwidth}L{0.25\textwidth}L{0.12\textwidth}L{0.26\textwidth}}
			Grado: & \answerline{10}\rule{0.8\linewidth}{0.4pt} & Fecha: & \rule{0.8\linewidth}{0.4pt} \\
			Nombre 1: & \multicolumn{3}{l}{\rule{0.63\textwidth}{0.4pt}} \\
			Nombre 2: & \multicolumn{3}{l}{\rule{0.63\textwidth}{0.4pt}} \\
		\end{tabular}
		
	\end{center}
	
	\section{Identificar a los usuarios}
	
	Revisen la observación realizada en la Clase 2. Identifiquen a las personas que experimentan o se ven afectadas por la situación.
	
	\vspace{-2.5mm}\begin{center}
		\small
		\begin{tabular}{|L{0.28\textwidth}|L{0.55\textwidth}|}
			\hline
			\textbf{Pregunta} & \textbf{Nuestra respuesta} \\
			\hline
			¿Qué situación estamos investigando? & \rule{0pt}{0.7cm}\begin{minipage}[b]{\linewidth}\sol{La espera para comprar alimentos en el punto de venta de la cafetería observado en C2.}\end{minipage} \\
			\hline
			¿Quiénes la experimentan o se ven afectados por ella? & \rule{0pt}{0.7cm}\begin{minipage}[b]{\linewidth}\sol{Estudiantes que compran o abandonan la fila y personal que atiende ese punto.}\end{minipage} \\
			\hline
			¿Por qué deberíamos preguntarles a estas personas? & \rule{0pt}{0.7cm}\begin{minipage}[b]{\linewidth}\sol{Conocen el efecto de la espera y condiciones del servicio que no podemos inferir al observar.}\end{minipage} \\
			\hline
		\end{tabular}
	\end{center}\vspace{-2.5mm}
	
	\section{Diseñar la encuesta}
	
	Creen una encuesta breve para conocer las experiencias, necesidades y dificultades de los usuarios. Las preguntas deben ser claras, neutrales y estar enfocadas en la situación.
	
	\vspace{-2.5mm}\begin{center}
		\small
		\begin{tabular}{|L{0.28\textwidth}|L{0.55\textwidth}|}
			\hline
			\textbf{Elemento} & \textbf{Nuestro plan de encuesta} \\
			\hline
			¿Quién responderá la encuesta? & \rule{0pt}{0.6cm}\begin{minipage}[b]{\linewidth}\sol{Estudiantes con experiencias distintas de compra; personal, en conversación complementaria.}\end{minipage} \\
			\hline
			¿Qué queremos aprender? & \rule{0pt}{0.7cm}\begin{minipage}[b]{\linewidth}\sol{Experiencias recientes, tiempo disponible, información para elegir y restricciones de atención.}\end{minipage} \\
			\hline
			¿Cómo recopilaremos las respuestas? & \rule{0pt}{0.6cm}\begin{minipage}[b]{\linewidth}\sol{Encuesta voluntaria y anónima fuera de la fila; aquí se modelan tres perfiles estudiantiles y uno de atención.}\end{minipage} \\
			\hline
		\end{tabular}
	\end{center}\vspace{-2.5mm}
	
	\vspace{0.15cm}
	
	\textbf{Escriban sus preguntas:}
	
	\vspace{-2.5mm}\begin{center}
		\small
		\begin{tabular}{|L{0.10\textwidth}|L{0.73\textwidth}|}
			\hline
			\textbf{\#} & \textbf{Pregunta de la encuesta} \\
			\hline
			1 & \rule{0pt}{0.65cm}\begin{minipage}[b]{\linewidth}\sol{En tu último descanso, ¿qué hiciste para conseguir alimentos y en qué momento?}\end{minipage} \\
			\hline
			2 & \rule{0pt}{0.65cm}\begin{minipage}[b]{\linewidth}\sol{Si esperaste para comprar, ¿cuánto calculas que esperaste antes de la atención?}\end{minipage} \\
			\hline
			3 & \rule{0pt}{0.65cm}\begin{minipage}[b]{\linewidth}\sol{¿Cómo decidiste qué comprar y qué información consultaste?}\end{minipage} \\
			\hline
			4 & \rule{0pt}{0.65cm}\begin{minipage}[b]{\linewidth}\sol{¿Qué ocurrió desde que llegaste a la fila hasta que recibiste tu compra o te retiraste?}\end{minipage} \\
			\hline
			5 & \rule{0pt}{0.65cm}\begin{minipage}[b]{\linewidth}\sol{Después de comprar o retirarte, ¿qué pudiste hacer en el tiempo restante del descanso?}\end{minipage} \\
			\hline
			6 & \rule{0pt}{0.65cm}\begin{minipage}[b]{\linewidth}\sol{¿Cómo supiste dónde hacer la fila y cuándo era tu turno?}\end{minipage} \\
			\hline
			7 & \rule{0pt}{0.65cm}\begin{minipage}[b]{\linewidth}\sol{En esa experiencia, ¿qué funcionó bien y qué dificultades, si hubo alguna, encontraste?}\end{minipage} \\
			\hline
		\end{tabular}
	\end{center}\vspace{-2.5mm}
	
	\newpage
	
	\textbf{Buenas preguntas para una encuesta:}
	
	\noindent\begin{tabular}{@{}p{0.4\linewidth}@{\hspace{0.04\linewidth}}p{0.55\linewidth}@{}}
		\begin{itemize}[leftmargin=0.6cm]
			\item Preguntan sobre experiencias reales;
			\item Son neutrales y claras;
			\item No sugieren una respuesta;
		\end{itemize}
		&
		\begin{itemize}[leftmargin=0.6cm]
			\item Ayudan a identificar necesidades o dificultades;
			\item Evitan pedir a los usuarios que elijan una solución.
		\end{itemize}
	\end{tabular}
	
	\vspace{3mm}
	
	\textbf{Ejemplo:} ``¿Con qué frecuencia experimenta esta situación?''
	
	\textbf{Eviten:} ``¿Esta situación es difícil o frustrante para usted?''
	
	\section{Realizar la encuesta}
	
	Realicen la encuesta con los usuarios seleccionados. Registren sus respuestas con la mayor precisión posible.
	
	\vspace{-2.5mm}\begin{center}
		\small
		\begin{tabular}{|L{0.22\textwidth}|L{0.73\textwidth}|}
			\hline
			\textbf{Elemento} & \textbf{Evidencia de los usuarios} \\
			\hline
			Usuario / perfil & \rule{0pt}{0.5cm}\begin{minipage}[b]{\linewidth}\sol{A: compra al inicio; B: compra más tarde; C: se retira; P: atiende. Perfiles anónimos simulados, no muestra representativa.}\end{minipage} \\
			\hline
			Experiencia relevante & \rule{0pt}{0.7cm}\begin{minipage}[b]{\linewidth}\sol{A (preguntas 1--4): llegó al inicio, calculó 7 min, preguntó precio al atenderlo y recibió su compra. B: llegó después, calculó 3 min y ya sabía qué pedir. C: se retiró al percibir poco tiempo. P: relata que debe cobrar, entregar y reponer.}\end{minipage} \\
			\hline
			Afirmación importante & \rule{0pt}{0.9cm}\begin{minipage}[b]{\linewidth}\sol{A (preguntas 5--7): ``Comí de afán; sabía mi turno, pero no el precio. B: ``Pude comer tranquilo; la fila avanzó. C: ``No sabía cuánto faltaba. P (conversación): ``Reponer también es parte de atender; hay que cobrar con cuidado.}\end{minipage} \\
			\hline
			Necesidad o frustración & \rule{0pt}{0.9cm}\begin{minipage}[b]{\linewidth}\sol{A necesita tiempo para comer e información para decidir; C, estimar si alcanza a comprar. B no reporta dificultad. P necesita atender con precisión y mantener productos disponibles.}\end{minipage} \\
			\hline
		\end{tabular}
	\end{center}\vspace{-2.5mm}
	
	\textbf{Importante:} Registren lo que los usuarios realmente dicen. No cambien sus respuestas para que coincidan con su idea original.
	
	\section{Comparar la evidencia}
	
	Comparen lo que dijeron los usuarios con lo que observaron en la Clase 2.
	%
	\vspace{-2.5mm}\begin{center}
		\small
		\begin{tabular}{|L{0.3\textwidth}|L{0.3\textwidth}|L{0.3\textwidth}|}
			\hline
			\textbf{Nuestra observación} & \textbf{Lo que dijeron los usuarios} & \textbf{Lo que aprendimos} \\
			\hline
			\rule{0pt}{1.3cm}\begin{minipage}[b]{\linewidth}\sol{En C2: 5--7 min en tres seguimientos; la fila baja después del minuto 6.}\end{minipage} & \rule{0pt}{1.3cm}\begin{minipage}[b]{\linewidth}\sol{A estima 7 min; B, 3 min al llegar más tarde.}\end{minipage} & \rule{0pt}{1.3cm}\begin{minipage}[b]{\linewidth}\sol{Coincidencia con A; B amplía la variación. Sus 3 min no fueron cronometrados en C2.}\end{minipage} \\
			\hline
			\rule{0pt}{1.3cm}\begin{minipage}[b]{\linewidth}\sol{Dos consultas de precios al iniciar atención; dos cruces junto a la fila.}\end{minipage} & \rule{0pt}{1.3cm}\begin{minipage}[b]{\linewidth}\sol{A consultó precio, pero reconoció su turno.}\end{minipage} & \rule{0pt}{1.3cm}\begin{minipage}[b]{\linewidth}\sol{Coincide la consulta; los cruces no prueban desorden ni pérdida de turnos.}\end{minipage} \\
			\hline
			\rule{0pt}{1.3cm}\begin{minipage}[b]{\linewidth}\sol{Atención activa y una pausa de reposición de 30 s.}\end{minipage} & \rule{0pt}{1.3cm}\begin{minipage}[b]{\linewidth}\sol{P describe reposición necesaria; C se retiró por incertidumbre.}\end{minipage} & \rule{0pt}{1.3cm}\begin{minipage}[b]{\linewidth}\sol{P cuestiona que toda pausa sea inútil; C aporta un abandono no registrado en C2.}\end{minipage} \\
			\hline
		\end{tabular}
	\end{center}\vspace{-2.5mm}
	
	Busquen evidencia que confirme, amplíe o contradiga sus observaciones.
	
	\section{¿Necesidad o solución?}
	
	Una necesidad está respaldada por evidencia sobre lo que las personas experimentan, necesitan o valoran. Una solución es algo que ustedes proponen.
	
	\vspace{-2.5mm}\begin{center}
		\small
		\begin{tabular}{|L{0.40\textwidth}|L{0.43\textwidth}|}
			\hline
			\textbf{Necesidad basada en evidencia} & \textbf{Posible solución} \\
			\hline
			Los usuarios necesitan pasar menos tiempo esperando. & Crear un sistema de reserva para autoservicio. \\
			\hline
			Los usuarios tienen dificultades para encontrar información. & Crear un sitio web con función de búsqueda. \\
			\hline
		\end{tabular}
	\end{center}\vspace{-2.5mm}
	
	\vspace{-2.5mm}\begin{center}
		\small
		\begin{tabular}{|L{0.30\textwidth}|L{0.53\textwidth}|}
			\hline
			\textbf{Pregunta} & \textbf{Nuestro análisis} \\
			\hline
			¿Qué necesidad está respaldada por evidencia? & \rule{0pt}{1.0cm}\begin{minipage}[b]{\linewidth}\sol{Comprar con tiempo suficiente para comer y descansar, y comprender la espera y el turno.}\end{minipage} \\
			\hline
			¿Qué solución estamos asumiendo? & \rule{0pt}{1.0cm}\begin{minipage}[b]{\linewidth}\sol{Podríamos anticipar una aplicación o más cajas; ninguna está justificada todavía.}\end{minipage} \\
			\hline
			¿Qué evidencia respalda la necesidad? & \rule{0pt}{1.0cm}\begin{minipage}[b]{\linewidth}\sol{Espera medida de 5--7 min; A come de afán y C se retira. B muestra experiencias distintas.}\end{minipage} \\
			\hline
		\end{tabular}
	\end{center}\vspace{-2.5mm}
	
	\vspace{-0.39cm}
	
	\section{Revisar nuestra comprensión}
	
	\vspace{-0.09cm}
	Utilicen la evidencia de las observaciones y de los usuarios para decidir si su comprensión original debe cambiar.
	
	\vspace{-2.5mm}\begin{center}
		\small
		\begin{tabular}{|L{0.27\textwidth}|L{0.27\textwidth}|L{0.27\textwidth}|}
			\hline
			\textbf{Idea original} & \textbf{Evidencia} & \textbf{¿Qué debería cambiar?} \\
			\hline
			\rule{0pt}{1.4cm}\begin{minipage}[b]{\linewidth}\sol{Las filas parecen excesivas e innecesarias.}\end{minipage} & \rule{0pt}{1.4cm}\begin{minipage}[b]{\linewidth}\sol{Hay atención continua, reposición y disminución de la fila.}\end{minipage} & \rule{0pt}{1.4cm}\begin{minipage}[b]{\linewidth}\sol{Distinguir espera evitable de tiempos necesarios; falta medir capacidad.}\end{minipage} \\
			\hline
			\rule{0pt}{1.4cm}\begin{minipage}[b]{\linewidth}\sol{La fila afecta por igual a todos.}\end{minipage} & \rule{0pt}{1.4cm}\begin{minipage}[b]{\linewidth}\sol{A come de afán; B come tranquilo; C abandona la compra.}\end{minipage} & \rule{0pt}{1.4cm}\begin{minipage}[b]{\linewidth}\sol{Precisar momentos y necesidades diferentes; no generalizar desde estos ejemplos.}\end{minipage} \\
			\hline
		\end{tabular}
	\end{center}\vspace{-2.5mm}
	
	\textbf{Nuestra conclusión:}
	
	\begin{multicols}{2}
		\begin{itemize}[leftmargin=0.6cm]
			\item \checkbox\quad Nuestra comprensión original está respaldada.
			\item \checkbox\makebox[0pt][r]{\sol{$\times$}\hspace{0.15em}}\quad Nuestra comprensión necesita ser modificada.
			\columnbreak
			\item \checkbox\quad Nuestra comprensión ha cambiado.
			\item \checkbox\makebox[0pt][r]{\sol{$\times$}\hspace{0.15em}}\quad Necesitamos más información.
		\end{itemize}
	\end{multicols}
	\vspace{-5mm}
	
	\textbf{Evidencia más sólida:}
	
	\noindent\noindent\answerline{Los tiempos de C2 y el relato de A coinciden; B muestra variación y P explica tareas.}\rule{\linewidth}{0.4pt}
	
	\vspace{0.25cm}
		\noindent\noindent\answerline{El caso es simulado y limitado: hacen falta más observaciones y voces de usuarios.}\rule{\linewidth}{0.4pt}
	
	\vspace{0.25cm}
	
	\textbf{Necesidad más importante de los usuarios:}
	
	\noindent\noindent\answerline{Acceder a alimentos con espera comprensible y margen para comer y descansar,}\rule{\linewidth}{0.4pt}
	
	\vspace{0.25cm}
		\noindent\noindent\answerline{con turnos claros y atención precisa; falta establecer qué demoras son evitables.}\rule{\linewidth}{0.4pt}
	
	\vspace{0.25cm}
	
	\textbf{Algo que aprendimos y que no esperábamos:}
	
	\noindent\noindent\answerline{Una fila puede avanzar y generar incertidumbre para quien considera comprar.}\rule{\linewidth}{0.4pt}
	
	\vspace{0.25cm}
	
	\noindent\noindent\answerline{La necesidad incluye prever el tiempo, además de reducir demoras evitables.}\rule{\linewidth}{0.4pt}
	
	\vspace{-0.29cm}
	
	\section{Verificación final}
	
	\noindent\begin{tabular}{@{}p{0.6\linewidth}@{\hspace{0.04\linewidth}}p{0.31\linewidth}@{}}
		\noindent\checkbox\makebox[0pt][r]{\sol{\(\times\)}\hspace{0.15em}}\quad Identificamos a los usuarios relevantes.
		
		\noindent\checkbox\makebox[0pt][r]{\sol{\(\times\)}\hspace{0.15em}}\quad Hicimos preguntas abiertas y registramos las respuestas.
		
		\vspace{0.12cm}
		
		\noindent\checkbox\makebox[0pt][r]{\sol{\(\times\)}\hspace{0.15em}}\quad Comparamos la evidencia con nuestras observaciones.
		
		\vspace{0.12cm}
		
		\noindent\checkbox\makebox[0pt][r]{\sol{\(\times\)}\hspace{0.15em}}\quad Identificamos una necesidad basada en evidencia.
		
		\vspace{0.12cm}
		
		\noindent\checkbox\makebox[0pt][r]{\sol{\(\times\)}\hspace{0.15em}}\quad Separamos las necesidades de las posibles soluciones.
		
		\vspace{0.12cm}
		
		\noindent\checkbox\makebox[0pt][r]{\sol{\(\times\)}\hspace{0.15em}}\quad Utilizamos evidencia para revisar nuestra comprensión original.
		&
		\vspace{-2.5mm}\begin{center}
			\textbf{Escuchen a los usuarios. Utilicen evidencia. Identifiquen la necesidad antes de proponer una solución.}
		\end{center}\vspace{-2.5mm}
	\end{tabular}
	
\end{document}
